\subsection{\texorpdfstring{\(\mathbf{A}_l\) 型代数 ( \(l \geq 1\) )}{\textbackslash mathbf\{A\}\_l 型代数 ( l \textbackslash geq 1 )}}

\begin{enumerate}
\def\labelenumi{(\Roman{enumi})}
\tightlist
\item
  令 V 为 k 上维数为 \(l+1\) 的向量空间，并令 \(\mathfrak{g}\) 为 V 的迹为零的自同态代数 \(\mathfrak{sl}(V)\)。令 \((e_{i})_{1\leq i\leq l+1}\) 为 V 的一组基；将 \(\mathfrak{g}\) 的元素关联到其关于这组基的矩阵的映射，是 \(\mathfrak{g}\) 与迹为零矩阵代数 \(\mathfrak{sl}(l+1,k)\) 之间的一个识别。我们知道 \(\mathfrak{g}\) 是半单的（第 I 章，第 6 节，第 7 号，命题 8）。
\end{enumerate}

回忆（代数，第 II 章，第 10 节，第 3 号），\(E_{ij}\) 表示矩阵 \((\alpha_{mp})\)，其中 \(\alpha_{ij} = 1\)，且 \(\alpha_{mp} = 0\) 对 \((m,p) \neq (i,j)\) 成立。矩阵

\[
\begin{array}{l l} E _ {i j} & (1 \leq i, j \leq l + 1, i \neq j) \\ E _ {i, i} - E _ {i + 1, i + 1} & (1 \leq i \leq l) \end{array}
\]

构成 \(\mathfrak{g}\) 的一组基。因此

\[
\dim \mathfrak {g} = l (l + 2).
\]

令 \(\hat{\mathfrak{h}}\) 为 \(\mathfrak{gl}(l + 1,k)\) 的对角元素集合；序列 \((E_{ii})_{1\leq i\leq l + 1}\) 是向量空间 \(\hat{\mathfrak{h}}\) 的一组基；令 \((\hat{\varepsilon}_i)_{1\leq i\leq l + 1}\) 为 \(\hat{\mathfrak{h}}^*\) 的一组基，且与 \((E_{ii})_{1\leq i\leq l + 1}\) 对偶。对所有 \(h\in \hat{\mathfrak{h}}\)，

\[
[ h, E _ {i j} ] = (\hat {\varepsilon} _ {i} (h) - \hat {\varepsilon} _ {j} (h)) E _ {i j}\tag{1}
\]

根据第 I 章，第 1 节，第 2 号的公式 (5)。令 \(\mathfrak{h}\) 为 \(\hat{\mathfrak{h}}\) 中迹为零的元素集合，并置 \(\varepsilon_i = \hat{\varepsilon}_i|\mathfrak{h}\)。则 \(\mathfrak{h}\) 是 \(\mathfrak{g}\) 的嘉当子代数（第 VII 章，第 2 节，第 1 号，例 4）。关系 (1) 证明此嘉当子代数是可分裂的，并且 \((\mathfrak{g},\mathfrak{h})\) 的根为 \(\varepsilon_i - \varepsilon_j\)（\(i\neq j\)）。令 \(\hat{\mathfrak{h}}_0^*\) 为 \(\hat{\mathfrak{h}}^*\) 中关于 \((\hat{\varepsilon}_i)\) 的坐标之和为零的元素集合。映射 \(\lambda \mapsto \lambda |\mathfrak{h}\) 从 \(\hat{\mathfrak{h}}_0^*\) 到 \(\mathfrak{h}^*\) 是双射。因此，\((\mathfrak{g},\mathfrak{h})\) 的根系 R 是 \(A_l\) 型（第 VI 章，第 4 节，第 7 号）。所以 \(\mathfrak{g}\) 是单的（第 3 节，第 2 号，命题 6 的推论 1）。因此，\(\mathfrak{g}\) 是 \(A_l\) 型的可分裂单李代数。

每个可分裂嘉当子代数 \(\mathfrak{h}^{\prime}\) 都是 \(\mathfrak{g}\) 中的 \(\mathfrak{h}\) 在某个初等自同构下的变换（\(\S3\)，第 3 号，命题 10 的推论）。由于 \(\operatorname{Aut}_{e}\mathfrak{g}\) 是由 \(x \mapsto sxs^{-1}\) 给出的 \(\mathfrak{g}\) 的自同构集合，其中 \(s \in \mathbf{SL}(V)\)（\(\S3\)，第 1 号，备注 2；另见 (VII)），存在 V 的一组基 \(\beta\)，使得 \(\mathfrak{h}^{\prime}\) 等于 \(\mathfrak{h}_{\beta}\)，即由 \(\mathfrak{g}\) 中关于基 \(\beta\) 的矩阵为对角矩阵的元素组成的集合。由于 \(\mathfrak{h}_{\beta}\) 含有一个具有互异特征值的元素，V 中在 \(\mathfrak{h}_{\beta}\) 的元素作用下稳定的向量子空间，只有由 \(\beta\) 的子集生成的那些。由此，映射 \(\beta \mapsto \mathfrak{h}_{\beta}\) 通过取商，诱导出从 V 分解为 \(l + 1\) 个一维子空间的分解集合到 \(\mathfrak{g}\) 的可分裂嘉当子代数集合的双射。

\begin{enumerate}
\def\labelenumi{(\Roman{enumi})}
\setcounter{enumi}{1}
\tightlist
\item
  令 \(\alpha = \varepsilon_{i} - \varepsilon_{j}\)（\(i \neq j\)）为一个根。有 \(\mathfrak{g}^{\alpha} = kE_{ij}\)。由于
\end{enumerate}

\[
[ E _ {i j}, E _ {j i} ] = E _ {i i} - E _ {j j}
\]

并且由于 \(\alpha(E_{ii} - E_{jj}) = 2\)，有（第 2 节，第 2 号，定理 1 (ii)）

\[
H _ {\alpha} = E _ {i i} - E _ {j j}.
\]

\begin{enumerate}
\def\labelenumi{(\Roman{enumi})}
\setcounter{enumi}{2}
\tightlist
\item
  置 \(\alpha_{1} = \varepsilon_{1} - \varepsilon_{2},\alpha_{2} = \varepsilon_{2} - \varepsilon_{3},\ldots ,\alpha_{l} = \varepsilon_{l} - \varepsilon_{l + 1}\)。根据第 VI 章，第 4 节，第 7 号.I，\((\alpha_{1},\dots,\alpha_{l})\) 是 R 的一组基 B；相对于 B 的正根是 \(\varepsilon_i - \varepsilon_j\)，其中 \(i < j\)。相应的 Borel 子代数 \(\mathfrak{b}\) 是迹为零的上三角矩阵集合。
\end{enumerate}

V 中的一个旗是 V 的向量子空间的一个集合，它不同于 \(\{0\}\) 和 V，并按包含关系全序。按包含关系给 V 的旗集合赋序。极大旗是集合 \(\{W_{1},\ldots,W_{l}\}\)，其中 \(W_{i}\) 是 i 维向量子空间，且

\[
\mathrm{W} _ {1} \subset \dots \subset \mathrm{W} _ {l}.
\]

例如，若 \(V_{i}\) 表示由 \(e_{1},\ldots,e_{i}\) 生成的 V 的子空间，则 \(\{V_{1},\ldots,V_{l}\}\) 是一个极大旗。

立即可见，\(\mathfrak{b}\) 是 \(\mathfrak{g}\) 中保持极大旗 \(\{V_{1},\ldots,V_{l}\}\) 的各元素不变的元素集合。反之，由于 \(\mathfrak{b}\) 包含 \(\mathfrak{h}\) 以及矩阵 \(E_{ij}\)，其中 i\textless j，可见 \(V_{i}\) 是在 \(\mathfrak{b}\) 下稳定的非平凡向量子空间中仅有的那些。

现在令 \(\delta\) 为 V 中的一个极大旗。由前述可知，\(\mathfrak{b}_{\delta}\)，即 \(\mathfrak{g}\) 中保持 \(\delta\) 的所有元素不变的元素集合，是 \(\mathfrak{g}\) 的一个 Borel 子代数。由于 \(\mathfrak{g}\) 的每个 Borel 子代数都是 \(\mathfrak{b}\) 在某个初等自同构下的变换，可知映射 \(\delta \mapsto \mathfrak{b}_{\delta}\) 是从极大旗集合到 \(\mathfrak{g}\) 的 Borel 子代数集合的双射。

令 \(\beta\) 为 V 的一组基。根据 (I) 及前述，包含 \(\mathfrak{h}_{\beta}\) 的 Borel 子代数正是与极大旗相对应的那些子代数，其中每个元素都由 \(\beta\) 的某个子集生成。这些旗与 \(\beta\) 上的全序一一对应，具体如下：该全序 \(\omega\) 对应于 \(\beta\) 上的旗 \(\{\mathrm{W}_1,\dots ,\mathrm{W}_l\}\)，其中 \(\mathrm{W}_i\) 是由前 \(i\) 个属于 \(\beta\) 的元素生成的向量子空间，对应的全序为 \(\omega\)。由于 \((l + 1)!\) 个全序定义在 \(\beta\) 上，我们重新得到 \((l + 1)!\) 个 \((\mathfrak{sl}(\mathrm{V}),\mathfrak{h}_{\beta})\) 中的 Borel 子代数这一事实（\(\S 3\)，第 3 号，备注）。

令 \(\gamma\) 为 V 中的一个旗。由于 \(\gamma\) 包含于某个极大旗，\(\mathfrak{p}_{\gamma}\)，即 \(\mathfrak{g}\) 中保持 \(\gamma\) 的各元素不变的元素集合，是 \(\mathfrak{g}\) 的一个抛物子代数。我们证明，在 \(\mathfrak{p}_{\gamma}\) 下稳定的非平凡向量子空间只有 \(\gamma\) 的元素。为此可设 \(\gamma = \{V_{i_1}, \ldots, V_{i_q}\}\)，其中 \(1 \leq i_1 < \cdots < i_q \leq l\)。置 \(i_0 = 0, i_{q+1} = l + 1\)。非空区间

\[
\left[ i _ {0} + 1, i _ {1} \right], \left[ i _ {1} + 1, i _ {2} \right], \dots , \left[ i _ {q} + 1, i _ {q + 1} \right]
\]

构成 \(\{1, \ldots, l + 1\}\) 的一个划分，从而任意 \(l + 1\) 阶方阵都可写成分块矩阵 \((X_{ab})_{1 \leq a, b \leq q + 1}\)。代数 \(\mathfrak{p}_{\gamma}\) 就是集合 \(\mathfrak{p}_{i_1,\dots,i_q}\)，它由元素 \((X_{ab})_{1\leq a,b\leq q + 1}\) 中属于 \(\mathfrak{sl}(l + 1,k)\) 且满足 \(X_{ab} = 0\)（当 \(a > b\)）的那些组成。由于 \(\mathfrak{p}_{i_1,\dots,i_q}\supset \mathfrak{b}\)，在 \(\mathfrak{p}_{i_1,\dots,i_q}\) 下稳定的非平凡向量子空间必为某个 \(\mathrm{V}_i\)；若 \(i_k < i < i_{k + 1}\)，代数 \(\mathfrak{p}_{i_1,\dots,i_q}\) 包含 \(E_{i_{k + 1},i}\)，而 \(\mathrm{V}_i\) 不稳定，故断言成立。

因此，\(2^{l}\) 个包含在极大旗 \(\{V_{1},\ldots,V_{l}\}\) 中的旗产生 \(2^{l}\) 个不同的抛物子代数，且这些子代数包含 \(\mathfrak{b}\)；由于共有 \(2^{l}\) 个包含 \(\mathfrak{b}\) 的抛物子代数（\(\S3\)，第 4 号，备注），可知映射 \(\gamma\mapsto \mathfrak{p}_{\gamma}\) 是从 V 的旗集合到 \(\mathfrak{g}\) 的抛物子代数集合的双射。此外，\(\mathfrak{p}_{\gamma}\supset \mathfrak{p}_{\gamma'}\) 当且仅当 \(\gamma\subset\gamma'\)。

回忆抛物子代数 \(\mathfrak{p} = \mathfrak{p}_{i_1, \ldots, i_q}\)（\(1 \leq i_1 < \cdots < i_q \leq l\)）。令 \(\mathfrak{s}\)（分别为 \(\mathfrak{n}\)）为由矩阵 \((X_{ab})_{1 \leq a, b \leq q+1}\) 构成的集合，这些矩阵属于 \(\mathfrak{sl}(l+1, k)\)，并满足 \(X_{ab} = 0\)：前者当 \(a \neq b\)，后者当 \(a \geq b\)。根据第 3 节第 4 号的命题 13，有 \(\mathfrak{p} = \mathfrak{s} \oplus \mathfrak{n}\)，子代数 \(\mathfrak{s}\) 在 \(\mathfrak{g}\) 中是约化的，而 \(\mathfrak{n}\) 既是 \(\mathfrak{p}\) 的最大幂零理想，也是其幂零根基。

\begin{enumerate}
\def\labelenumi{(\Roman{enumi})}
\setcounter{enumi}{3}
\tightlist
\item
  对 \(r = 1, 2, \ldots, l\)，令 \(\varpi_{r} = \varepsilon_{1} + \cdots + \varepsilon_{r}\)。有 \(\varpi_{i}(H_{\alpha_{j}}) = \delta_{ij}\)，所以 \(\varpi_{r}\) 是对应于 \(\alpha_{r}\) 的基本权。
\end{enumerate}

令 \(\sigma\) 为 \(\mathfrak{g}\) 在 V 上的恒等表示。\(\bigwedge^{r}\sigma\) 是 \(\sigma\) 在 \(\mathrm{E} = \bigwedge^{r}(\mathrm{V})\) 上的表示。令 \((e_1,\ldots ,e_{l + 1})\) 为 V 的所选基。这些 \(e_{i_1}\wedge \dots \wedge e_{i_r}\)（其中 \(i_1 < \dots < i_r\)）构成 E 的一组基。若 \(h\in \mathfrak{h}\)，

\[
(\bigwedge^ {r} \sigma) (h). e _ {i _ {1}} \wedge \dots \wedge e _ {i _ {r}} = (\varepsilon_ {i _ {1}} + \dots + \varepsilon_ {i _ {r}}) (h) e _ {i _ {1}} \wedge \dots \wedge e _ {i _ {r}}.
\]

因此，每个权的重数都是 1，\(\varpi_{r}\) 是 \(\bigwedge^{r}\sigma\) 的一个权，而其他每个权都形如 \(\varpi_{r}-\mu\)，其中 \(\mu\) 是正根权。所以，\(\varpi_{r}\) 是 \(\bigwedge^{r}\sigma\) 的最高权，而 \(e_{1}\wedge\cdots\wedge e_{r}\) 是一个本原元。根据第 VI 章，第 4 节，第 7 号.IX，Weyl 群可识别为集合

\[
\{\varepsilon_ {1}, \dots , \varepsilon_ {l + 1} \}.
\]

因此，\(\varpi_{r}\) 在 Weyl 群下的轨道包含所有 \(\varepsilon_{i_{1}} + \cdots + \varepsilon_{i_{r}}\)，其中 \(i_{1} < \cdots < i_{r}\)。由本原元 \(e_{1} \wedge \cdots \wedge e_{r}\) 生成的单子模因而以所有 \(\varepsilon_{i_{1}} + \cdots + \varepsilon_{i_{r}}\) 为权，从而等于 E。因此，\(\bigwedge^{r} \sigma\) 是最高权为 \(\varpi_{r}\) 的不可约表示。

因此，表示 \(\bigwedge^{r}\sigma\)（\(1 \leq r \leq l\)）是基本表示。有 \(\dim(\bigwedge^{r}\sigma) = \binom{l+1}{r}\)。

\begin{enumerate}
\def\labelenumi{(\Alph{enumi})}
\setcounter{enumi}{21}
\tightlist
\item
  有 \(w_0(\alpha_1) = -\alpha_l, w_0(\alpha_2) = -\alpha_{l-1}, \ldots\)（第 VI 章，第 4 节，第 7 号，XI），因此
\end{enumerate}

\[
- w _ {0} (\varpi_ {1}) = \varpi_ {l}, - w _ {0} (\varpi_ {2}) = \varpi_ {l - 1}, \ldots .
\]

令

\[
\omega = n _ {1} \varpi_ {1} + \dots + n _ {l} \varpi_ {l} \quad (n _ {1}, \ldots , n _ {l} \in \mathbf {N})
\]

为一个支配权。具有最高权 \(\omega\) 的单表示是正交型或辛型，当且仅当

\[
n _ {1} = n _ {l}, \quad n _ {2} = n _ {l - 1}, \ldots
\]

（第 7 节，第 5 号，命题 12）。特别地，若 \(l\) 为偶数，则 \(\mathfrak{sl}(l + 1,k)\) 的基本表示没有一个是正交型或辛型。若 \(l\) 为奇数，则表示 \(\bigwedge^i\sigma\)（其中 \(i\neq (l + 1) / 2\)）既非正交型也非辛型；根据第 VI 章，第 4 节，第 7 号.VI，\(\varpi_{(l + 1) / 2}\) 关于 \((\alpha_{1},\ldots ,\alpha_{l})\) 的坐标之和为

\[
\begin{array}{r l} \frac {1}{l + 1} \left[ \frac {l + 1}{2} \left(1 + 2 + \dots + \frac {l - 1}{2}\right) + \frac {l + 1}{2} \left(1 + 2 + \dots + \frac {l + 1}{2}\right) \right] & = 1 + 2 + \dots + \frac {l - 1}{2} + \frac {l + 1}{4} \end{array}
\]

所以，\(\bigwedge^{(l + 1) / 2}\sigma\) 在 \(l\equiv -1\)（mod. 4）时是正交型的；在 \(l\equiv 1\)（mod. 4）时是辛型的（第 7 节，第 5 号，命题 12）。这一结果还可如下精确化。选取一个非零元素 \(e\)，属于 \(\bigwedge^{l + 1}(\mathrm{V})\)。外代数 \(\mathrm{V}\) 中的乘法定义了从

\[
\bigwedge^ {(l + 1) / 2} (\mathrm{V}) \times \bigwedge^ {(l + 1) / 2} (\mathrm{V})
\]

到 \(\bigwedge^{l + 1}(\mathrm{V})\) 的双线性映射，该映射可写成 \((u,v)\mapsto \varPhi (u,v)e\)，其中 \(\varPhi\) 是 \(\bigwedge^{(l + 1) / 2}(\mathrm{V})\) 上的双线性型。立即可验证，\(\varPhi\) 非零且在 \(\mathfrak{g}\) 下不变（因而非退化）；当 \((l + 1) / 2\) 为偶数时它对称，当 \((l + 1) / 2\) 为奇数时它交错。

\begin{enumerate}
\def\labelenumi{(\Roman{enumi})}
\setcounter{enumi}{5}
\tightlist
\item
  对所有 \(x \in \mathfrak{g}\)，\(\sigma(x) = x\) 的特征多项式可写成
\end{enumerate}

\[
\mathrm{T} ^ {l + 1} + f _ {2} (x) \mathrm{T} ^ {l - 1} + f _ {3} (x) \mathrm{T} ^ {l - 2} + \dots + f _ {l + 1} (x)
\]

其中 \(f_{2}, \ldots, f_{l+1}\) 是在 \(\mathfrak{g}\) 下不变的多项式函数（第 8 节，第 3 号，引理 2）。

若 \(x = \xi_{1}E_{11} + \cdots + \xi_{l+1}E_{l+1 l+1} \in \mathfrak{h}\)，则 \(f_{i}(x)\) 除了符号外，正是 \(\xi_{1}, \ldots, \xi_{l+1}\) 的次数为 \(2, \ldots, l + 1\) 的初等对称函数。因此，根据第 VI 章，第 4 节，第 7 号.IX，\(f_{i}|_{\mathfrak{h}}\) 生成 \(\mathbf{S}(\mathfrak{h}^{*})\) 中在 Weyl 群下不变的元素所成的代数，并且代数无关。所以（第 8 节，第 3 号，命题 3），\(f_{2}, f_{3}, \ldots, f_{l+1}\) 生成在 \(\mathfrak{g}\) 下不变的多项式函数所成的代数，并且代数无关。

\begin{enumerate}
\def\labelenumi{(\Roman{enumi})}
\setcounter{enumi}{6}
\tightlist
\item
  对所有 \(g \in \mathbf{GL}(l + 1, k)\)，令 \(\varphi_k(g) = \varphi(g)\) 为由 \(x \mapsto gxg^{-1}\) 给出的 \(\mathfrak{g}\) 的自同构。则 \(\varphi\) 是从 \(\mathbf{GL}(l + 1, k)\) 到 \(\operatorname{Aut}(\mathfrak{g})\) 的同态。有
\end{enumerate}

\[
\varphi (\mathbf {S L} (l + 1, k)) = \operatorname{Aut} _ {e} (\mathfrak {g})
\]

（第 VII 章，第 3 节，第 1 号，备注 2）。令 \(\bar{k}\) 为 \(k\) 的代数闭包。有

\[
\mathbf {G L} (l + 1, \bar {k}) = \bar {k} ^ {*}. \mathbf {S L} (l + 1, \bar {k}),
\]

所以 \(\varphi_{\bar{k}}(\mathbf{GL}(l + 1,\bar{k})) = \varphi_{\bar{k}}(\mathbf{SL}(l + 1,\bar{k})) = \mathrm{Aut}_e(\mathfrak{g}\otimes_k\bar{k})\)；由此 \(\varphi (\mathbf{GL}(l + 1,k))\subset \mathrm{Aut}_0(\mathfrak{g})\)。另一方面，\(\mathrm{Aut}_0(\mathfrak{g})\subset \varphi (\mathbf{GL}(l + 1,k))\) 由将第 7 节第 1 号的命题 2 应用于 \(\mathfrak{g}\) 的恒等表示而得。因此，

\[
\operatorname{Aut} _ {0} (\mathfrak {g}) = \varphi (\mathbf {G L} (l + 1, k)).
\]

 \(\varphi\) 的核是 \(\mathbf{GL}(l+1,k)\) 中与每个 \(l+1\) 阶矩阵对易的元素集合，即可逆标量矩阵的集合 \(k^{*}\)。因此，\(\operatorname{Aut}_{0}(\mathfrak{g})\) 可识别为群 \(\mathbf{GL}(l+1,k)/k^{*}=\mathbf{PGL}(l+1,k)\)。\(\varphi'=\varphi|\mathbf{SL}(l+1,k)\) 的核是 \(\mu_{l+1}(k)\)，其中 \(\mu_{l+1}(k)\) 表示 k 中的 \((l+1)\) 次单位根集合。因此，\(\operatorname{Aut}_{e}(\mathfrak{g})\) 可识别为群 \(\mathbf{SL}(l+1,k)/\mu_{l+1}(k)=\mathbf{PSL}(l+1,k)\)。另一方面，有正合序列

\[
1 \longrightarrow \mathbf {S L} (l + 1, k) \longrightarrow \mathbf {G L} (l + 1, k) \stackrel {\det} {\longrightarrow} k ^ {*} \longrightarrow 1
\]

且 \(k^*\) 在 det 下的像为 \(k^{*l + 1}\)。由此得到典则同构

\[
\begin{array}{c} \operatorname{Aut} _ {0} (\mathfrak {g}) / \operatorname{Aut} _ {e} (\mathfrak {g}) \longrightarrow \mathbf {P G L} (l + 1, k) / \mathbf {P S L} (l + 1, k) \\ \qquad \qquad \qquad \qquad \qquad \longrightarrow \mathbf {G L} (l + 1, k) / k ^ {*}. \mathbf {S L} (l + 1, k) \longrightarrow k ^ {*} / k ^ {*   l + 1}. \end{array}
\]

若 \(k = \mathbf{R}\)，可见 \(\mathrm{Aut}_0(\mathfrak{g}) = \mathrm{Aut}_e(\mathfrak{g})\) 当 \(l + 1\) 为奇数时成立，而 \(\mathrm{Aut}_0(\mathfrak{g}) / \mathrm{Aut}_e(\mathfrak{g})\) 同构于 \(\mathbf{Z}/2\mathbf{Z}\) 当 \(l + 1\) 为偶数时。

采用第 5 节的记号，\(f(\mathrm{T}_{\mathbf{Q}})\) 是 \(\mathfrak{g}\) 的在 \(\mathfrak{h}\) 上诱导恒等映射的自同构集合，因而等于 \(\varphi(\mathrm{D})\)，其中 D 是 \(\mathbf{GL}(l+1,k)\) 的对角元素集合（第 5 节，命题 4）。令 \(D'\) 为 \(\mathbf{SL}(l+1,k)\) 的对角元素集合。根据第 5 节的命题 3 以及 \(\operatorname{Aut}_{e}(\mathfrak{g})\) 的确定，有 \(f(q(\mathrm{T}_{\mathrm{P}})) \subset \varphi(\mathrm{D}')\)。我们证明 \(f(q(\mathrm{T}_{\mathrm{P}})) = \varphi(\mathrm{D}')\)。令

\[
d = \left( \begin{array}{c c c} \lambda_ {1} & & 0 \\ & \ddots & \\ 0 & & \lambda_ {l + 1} \end{array} \right)
\]

为 \(D'\) 的一个元素。存在 \(\zeta \in \operatorname{Hom}(\mathbf{Q}(\mathbf{R}), k^{*}) = \mathrm{T}_{\mathbf{Q}}\)，使得对所有 i 和 j 都有 \(\zeta(\varepsilon_{i} - \varepsilon_{j}) = \lambda_{i}\lambda_{j}^{-1}\)。容易验证 \(f(\zeta) = \varphi(d)\)。根据第 VI 章，第 4 节，第 7 号.VIII，\(\mathrm{P}(\mathrm{R})\) 由 \(\mathrm{Q}(\mathrm{R})\) 和元素 \(\varepsilon = \varepsilon_{1}\) 生成，后者在 \(\mathrm{P}(\mathrm{R})/\mathrm{Q}(\mathrm{R})\) 中的像的阶为 \(l + 1\)；但是

\[
\begin{array}{c} \zeta ((l + 1) \varepsilon) = \zeta ((\varepsilon_ {1} - \varepsilon_ {2}) + (\varepsilon_ {1} - \varepsilon_ {3}) + \dots + (\varepsilon_ {1} - \varepsilon_ {l + 1})) \\ = \lambda_ {1} ^ {l} \lambda_ {2} ^ {- 1} \lambda_ {3} ^ {- 1} \ldots \lambda_ {l + 1} ^ {- 1} = \lambda_ {1} ^ {l + 1} \end{array}
\]

所以 \(\zeta\) 延拓为从 \(\mathrm{P}(\mathrm{R})\) 到 \(k^{*}\) 的同态。这证明 \(\zeta \in q(\mathrm{T}_{\mathrm{P}})\)，从而 \(\varphi(d) \in f(q(\mathrm{T}_{\mathrm{P}}))\)。

回忆（第 5 节，第 3 号，命题 5 的推论 2），\(\mathrm{Aut}(\mathfrak{g}) = \mathrm{Aut}_0(\mathfrak{g})\) 当 \(l = 1\) 时，且 \(\mathrm{Aut}(\mathfrak{g}) / \mathrm{Aut}_0(\mathfrak{g})\) 同构于 \(\mathbf{Z}/2\mathbf{Z}\) 当 \(l \geq 2\) 时。映射 \(\theta : x \mapsto -^t x\) 是 \(\mathfrak{sl}(l + 1, k)\) 的一个自同构，且 \(a_0 = \theta|\mathfrak{h} \notin \mathrm{W}\) 当 \(l \geq 2\) 时（第 VI 章，第 4 节，第 7 号.XI），所以 \(a_0\) 在 \(\mathrm{Aut}(\mathfrak{g}) / \mathrm{Aut}_0(\mathfrak{g})\) 中的类是该群的非平凡元素（第 5 节，第 2 号，命题 4）。

\begin{enumerate}
\def\labelenumi{(\Roman{enumi})}
\setcounter{enumi}{7}
\tightlist
\item
  Killing 型在 \(\mathfrak{h}\) 上的限制为
\end{enumerate}

\[
\begin{array}{l} \Phi (\xi_ {1} E _ {1 1} + \dots + \xi_ {l + 1} E _ {l + 1, l + 1}, \xi_ {1} ^ {\prime} E _ {1 1} + \dots + \xi_ {l + 1} ^ {\prime} E _ {l + 1, l + 1}) \\ = \sum_ {i \neq j} (\xi_ {i} - \xi_ {j}) (\xi_ {i} ^ {\prime} - \xi_ {j} ^ {\prime}) = \sum_ {i, j} (\xi_ {i} - \xi_ {j}) (\xi_ {i} ^ {\prime} - \xi_ {j} ^ {\prime}) \\ = (l + 1) \sum_ {i} \xi_ {i} \xi_ {i} ^ {\prime} + (l + 1) \sum_ {j} \xi_ {j} \xi_ {j} ^ {\prime} - 2 \left(\sum_ {i} \xi_ {i}\right) \left(\sum_ {j} \xi_ {j} ^ {\prime}\right) \\ = 2 (l + 1) \sum_ {i} \xi_ {i} \xi_ {i} ^ {\prime}. \end{array}
\]

\begin{enumerate}
\def\labelenumi{(\Roman{enumi})}
\setcounter{enumi}{8}
\tightlist
\item
  对 \(1 \leq i < j \leq l + 1\)，置
\end{enumerate}

\[
X _ {\varepsilon_ {i} - \varepsilon_ {j}} = E _ {i j} X _ {\varepsilon_ {j} - \varepsilon_ {i}} = - E _ {j i}.
\]

那么，对所有 \(\alpha \in \mathrm{R}\)，有 \([X_{\alpha}, X_{-\alpha}] = -H_{\alpha}\) 和 \(\theta(X_{\alpha}) = X_{-\alpha}\)（其中 \(\theta\) 是在（VII）中引入的自同构 \(x \mapsto -^{t} x\)）。因此，\((X_{\alpha})_{\alpha \in \mathrm{R}}\) 是 \((\mathfrak{g}, \mathfrak{h})\) 中的一个谢瓦莱系。

取 \(k = \mathbf{Q}\)。\(\mathfrak{h}\) 中的许可格（第 12 节，第 6 号，定义 1）正是介于由 \(\mathbf{Z}\)-模 \(\mathrm{Q}(\mathrm{R}^{\vee})\) 生成的 \(E_{ii} - E_{i+1,i+1}\)（即由属于 \(\mathfrak{sl}(l+1,\mathbf{Z})\) 的对角矩阵组成）与由 \(\mathbf{Z}\)-模 \(\mathrm{P}(\mathrm{R}^{\vee})\) 生成的 \(\mathrm{Q}(\mathrm{R}^{\vee})\) 及 \(E_{11} - (l+1)^{-1}\sum E_{ii}\)（第 VI 章，第 4 节，第 7 号.VIII，即由形如 \(x + (l+1)^{-1}a.1\) 的迹为零对角矩阵组成，其中 \(x\) 的元素为整数且 \(a\in\mathbf{Z}\)）之间的那些格。因此，\(\mathfrak{sl}(l+1,\mathbf{Z})\) 是 \((\mathfrak{g},\mathfrak{h})\) 中与许可格 \(\mathrm{Q}(\mathrm{R}^{\vee})\) 及谢瓦莱系 \((X_{\alpha})\) 相联系的谢瓦莱序。容易验证，\(\bigwedge^{r}\mathbf{Z}^{l+1}\) 是 \(\bigwedge^{r}\mathbf{Q}^{l+1}\) 中相对于 \(\mathfrak{sl}(l+1,\mathbf{Z})\) 的一个容许格（第 12 节，第 8 号，定义 3）。

另一方面，\(\mathfrak{gl}(l+1,\mathbf{Z})\) 是分裂约化代数 \(\mathfrak{gl}(l+1,\mathbf{Q})\) 中的谢瓦莱序；它在 \(\mathfrak{sl}(l+1,\mathbf{Q})\) 上的投影平行于中心 \(\mathbf{Q}.1\)（\(\mathfrak{gl}(l+1,\mathbf{Q})\) 的中心），所得的是 \((\mathfrak{g},\mathfrak{h})\) 中由许可格 \(\mathrm{P}(\mathrm{R}^{\vee})\)（位于 \(\mathfrak{h}\) 中）和谢瓦莱系 \((X_{\alpha})\) 定义的谢瓦莱序。我们指出，\(\mathfrak{gl}(l+1,\mathbf{Z})\) 不是它与 \(\mathfrak{sl}(l+1,\mathbf{Q})\) 的交及其与 \(\mathfrak{gl}(l+1,\mathbf{Q})\) 的中心的交的直和。
% semantic-fragments: legacy-input-seam
\relax{}
